\chapter{Architecture et conception du système}

\section{Choix architectural : microservices vs monolithe}

\section{Vue globale de l'architecture}

% \begin{figure}[H]
%   \centering
%   \includegraphics[width=\textwidth]{architecture-globale}
%   \caption{Architecture globale d'ASM Track}
% \end{figure}

\section{Description des microservices}

\begin{table}[H]
\centering
\caption{Microservices d'ASM Track}
\begin{tabularx}{\textwidth}{|L{3.5cm}|C{1.5cm}|X|}
\hline
\rowcolor{primary!10}
\textbf{Service} & \textbf{Port} & \textbf{Responsabilité} \\
\hline
AuthServer & 8085 & Émission et validation des JWT (OAuth2) \\
\hline
ApiGateway & 8080 & Routage, authentification, RBAC \\
\hline
AppBackend & 8081 & Gestion des utilisateurs, entreprises, véhicules \\
\hline
DeliveryMicroservice & 8082 & Cycle de vie des livraisons, tournées, rapports \\
\hline
DriverService & 8083 & Profils chauffeurs, statistiques \\
\hline
ErpAdapterService & 8084 & Intégration ERP agnostique \\
\hline
\end{tabularx}
\end{table}

\section{Choix technologiques}

\subsection{Backend}

\subsection{Frontend}

\subsection{Mobile}

\subsection{Infrastructure}

\section{Schéma de déploiement Docker}

\section*{Conclusion}
